\section{Ideas and Future Work}
\label{sec:future}

\begin{enumerate}
  \item \textbf{Main Ambition: even better up-to with diacritical companions}

        Imagine a proof by coinduction of \emph{weak} bisimilarity (bisimilarity
        modulo invisible steps). We can rewrite by \textbf{strong bisimilarity},
        but \textbf{not by weak bisimilarity} in general. There are special cases
        where we can, though! How? Only if we use CIH under beta, not tau.
        \code{gpaco} lets us do this, and so does the \emph{diacritical
        companion}~\cite{diacritical}, a newer discovery. Can we recover this power in the land of
        the tower?

        There exists modern work which could recover some of the power of gpaco
        within the tower framework, but if this is usable remains to be seen
        (part of my summer work at Inria).

  \item \textbf{A write-up on the tower in the library}

        Because history is nonlinear, no write-up of the enhanced coinduction
        library as it is exists.

  \item \textbf{Immediate: Upstream to ITrees/rocq-coinduction}

        Usability of the enhanced coinduction library is good, but there are
        bugs. Also, there are typical design patterns with the library that
        probably should be part of it.
\end{enumerate}

\note{Item 2 is a contribution the paper can partly claim rather than defer: the
rewritten library ships with its test suite grown $84 \to 328$ lines and its
comments $370 \to 455$. Item 3's ``design patterns that should be part of it'' now
has a worked instance, the shared \code{to_mon} driver, which removed the
per-functor blocks from three downstream files and took \code{extra/}'s tactic
cost from $895$ characters below paco's $405$, to $262$.}
